\documentclass[11pt,a4paper]{article}
\usepackage[utf8]{inputenc}
\usepackage[T1]{fontenc}
\usepackage[spanish]{babel}
\usepackage{amsmath,amssymb}
\usepackage{booktabs}
\usepackage[margin=2.5cm]{geometry}
\usepackage{hyperref}


\newcommand{\Ndet}{\textbf{[6]}}        % numero de estrellas detectadas
\newcommand{\Nnodet}{\textbf{[94]}}    % numero de estrellas no detectadas
\newcommand{\Pdet}{\textbf{[6]}}        % porcentaje de detectadas
\newcommand{\Pnodet}{\textbf{[94]}}        % porcentaje de detectadas
\newcommand{\Fdet}{\textbf{[3.896869e-10]}}        % flujo promedio, detectadas
\newcommand{\Fnodet}{\textbf{[1.794581e-10]}}    % flujo promedio, no detectadas
\newcommand{\ddet}{\textbf{[76.682944]}}        % distancia promedio, detectadas
\newcommand{\dnodet}{\textbf{[274.024742]}}    % distancia promedio, no detectadas
\newcommand{\Tdet}{\textbf{[9049.768732]}}        % temperatura promedio, detectadas
\newcommand{\Tnodet}{\textbf{[6517.771301]}}    % temperatura promedio, no detectadas
\newcommand{\Lmax}{\textbf{[162.5988]}}        % luminosidad maxima (L_sol)
\newcommand{\Lmin}{\textbf{[0.0597]}}        % luminosidad minima (L_sol)

\title{Simulación de un catálogo estelar: luminosidad, flujo observado y detección}
\author{Valentín Carroza\\ \small  Física Computacional}
\date{\today}

\begin{document}
\maketitle
% ==================================================================
\section{Introducción}
% ==================================================================
El objetivo de este trabajo es simular un catálogo astronómico y estudiar qué estrellas
puede detectar un instrumento con una sensibilidad dada. Para ello se relacionan las propiedades
intrínsecas de una estrella (radio y temperatura efectiva) con la cantidad de energía que
llega a un observador, que depende además de la distancia.

La energía total que una estrella emite por unidad de tiempo, o luminosidad, se obtiene al
tratarla como un cuerpo negro de radio $R$ y temperatura efectiva $T_{\mathrm{eff}}$. La ley de
Stefan-Boltzmann establece que
\begin{equation}
  L = 4\pi R^{2}\,\sigma_{\mathrm{SB}}\,T_{\mathrm{eff}}^{4},
  \label{eq:stefan}
\end{equation}
donde $\sigma_{\mathrm{SB}} = 5{,}670\times10^{-8}\ \mathrm{W\,m^{-2}\,K^{-4}}$ es la constante de Stefan-Boltzmann.

El flujo observado es la energía que atraviesa una unidad de área por unidad de tiempo a una distancia $d$ de la fuente:
\begin{equation}
  F = \frac{L}{4\pi d^{2}}.
  \label{eq:flujo}
\end{equation}
Esta cantidad es la que mide realmente un telescopio. Dos estrellas con la misma luminosidad
pueden tener flujos muy distintos si están a distancias diferentes, por lo que el flujo
determina, junto con la sensibilidad del instrumento, qué objetos pueden ser detectados.

% ==================================================================
\section{Metodología}
% ==================================================================
\subsection{Generación del catálogo}
Se generó una muestra de $N=100$ estrellas con \texttt{NumPy}, utilizando un generador
pseudoaleatorio con semilla fija para que los resultados sean reproducibles. Las tres
propiedades se obtuvieron de distribuciones uniformes independientes:
\begin{itemize}
  \item radio estelar: $R \sim \mathcal{U}(0{,}5,\,5)\,R_\odot$;
  \item distancia: $d \sim \mathcal{U}(50,\,500)\ \mathrm{pc}$;
  \item temperatura efectiva: $T_{\mathrm{eff}} \sim \mathcal{U}(3000,\,10000)\ \mathrm{K}$.
\end{itemize}
Los datos se almacenaron en un \texttt{DataFrame} de \texttt{Pandas} con un identificador por estrella y una columna por propiedad.

\subsection{Cálculo de la luminosidad y del flujo}
Las propiedades se convirtieron al Sistema Internacional ($R_\odot = 6{,}957\times10^{8}$~m,
$1\ \mathrm{pc} = 3{,}0857\times10^{16}$~m). La luminosidad se calculó con la
ecuación~\eqref{eq:stefan} y se expresó en luminosidades solares dividiendo por
$L_\odot = 3{,}828\times10^{26}$~W. El flujo se calculó con la ecuación~\eqref{eq:flujo},
usando $L$ en watts y $d$ en metros, de modo que el resultado queda en $\mathrm{W\,m^{-2}}$.

\subsection{Criterio de clasificación}
Se supuso un telescopio capaz de detectar estrellas con flujo igual o superior a
$F_{\mathrm{lim}} = 2\times10^{-10}\ \mathrm{W\,m^{-2}}$. Cada estrella se clasificó como
\emph{Detectada} si $F \geq F_{\mathrm{lim}}$ y como \emph{No detectada} si
$F < F_{\mathrm{lim}}$.

% ==================================================================
\section{Resultados}
% ==================================================================
La luminosidad del catálogo varía entre $\Lmin\,L_\odot$ y $\Lmax\,L_\odot$.

El resultado de la clasificación se resume en la Tabla~\ref{tab:clasificacion}.

\begin{table}[ht]
  \centering
  \caption{Resultados principales de la clasificación para $F_{\mathrm{lim}} = 2\times10^{-10}\
  \mathrm{W\,m^{-2}}$. Los promedios se calcularon por separado para cada categoría.}
  \label{tab:clasificacion}
  \begin{tabular}{lccccc}
    \toprule
    Categoría & $N$ & Porcentaje (\%) & $\langle F\rangle$ ($\mathrm{W\,m^{-2}}$) &
    $\langle d\rangle$ (pc) & $\langle T_{\mathrm{eff}}\rangle$ (K) \\
    \midrule
    Detectada    & \Ndet   & \Pdet & \Fdet   & \ddet   & \Tdet   \\
    No detectada & \Nnodet & \Pnodet & \Fnodet & \dnodet & \Tnodet \\
    \bottomrule
  \end{tabular}
\end{table}

% ==================================================================
\section{Discusión}
% ==================================================================
\paragraph{Dependencia de la luminosidad con la temperatura y el radio.}
Según la ecuación~\eqref{eq:stefan}, $L \propto R^{2}T_{\mathrm{eff}}^{4}$. La dependencia
con la temperatura es mucho más fuerte: duplicar $T_{\mathrm{eff}}$ multiplica la luminosidad por 16, mientras que duplicar $R$ la multiplica solo por 4. Esto se observa como un aumento de la luminosidad hacia las temperaturas altas, con una dispersión
vertical causada por las diferencias de radio.

\paragraph{Relación entre distancia y flujo.} 
De la ecuación~\eqref{eq:flujo}, $F \propto d^{-2}$ para una luminosidad dada. Se aprecia una tendencia decreciente del flujo con la distancia, pero con una dispersión considerable, porque el flujo también depende de $R$ y $T_{\mathrm{eff}}$.

\paragraph{¿Las estrellas cercanas son necesariamente detectadas?}
No. La detección depende del flujo, que combina la luminosidad intrínseca y la distancia. Una estrella cercana pero poco luminosa (pequeña y fría) puede tener un flujo inferior a
$F_{\mathrm{lim}}$, mientras que una estrella más lejana pero muy luminosa puede superarlo.

\paragraph{Influencia de la temperatura y del radio en la detección.}
Como $F = \sigma_{\mathrm{SB}}\,T_{\mathrm{eff}}^{4}\,(R/d)^{2}$, la condición de detección
$F \geq F_{\mathrm{lim}}$ equivale a $T_{\mathrm{eff}}^{2}\,R/d \geq \sqrt{F_{\mathrm{lim}}/\sigma_{\mathrm{SB}}}$ (con $R$ y $d$ en las mismas unidades). La temperatura entra al cuadrado y el radio solo de forma lineal, de modo que las estrellas calientes tienen una ventaja mayor para ser detectadas. Un radio grande ayuda, pero no compensa una temperatura baja ni una distancia grande.

\paragraph{Efecto de disminuir el límite de detección.}
Si $F_{\mathrm{lim}}$ disminuyera, el telescopio sería más sensible y la condición $F \geq F_{\mathrm{lim}}$ la cumplirían más estrellas; el número de estrellas detectadas aumentaría o, como mínimo, no disminuiría.

\paragraph{Estrellas con la misma temperatura y distinta luminosidad.}
Porque la luminosidad depende también del radio: $L \propto R^{2}T_{\mathrm{eff}}^{4}$. A igual temperatura, la estrella de mayor radio tiene más superficie emisora y, por lo tanto, mayor luminosidad.

% ==================================================================
\section{Conclusiones}
% ==================================================================
Se generó un catálogo de 100 estrellas y se calcularon su luminosidad y su flujo observado
mediante la ley de Stefan-Boltzmann. De los resultados se concluye que:
\begin{itemize}
  \item La luminosidad crece como $R^{2}T_{\mathrm{eff}}^{4}$ y está dominada por la temperatura;
  \item El flujo decrece aproximadamente como $d^{-2}$, aunque la dispersión es grande porque depende también de $R$ y $T_{\mathrm{eff}}$;
  \item Con $F_{\mathrm{lim}} = 2\times10^{-10}\ \mathrm{W\,m^{-2}}$ se detectaron \Ndet\ de las 100 estrellas (\Pdet\,\%);
  \item La cercanía por sí sola no garantiza la detección, y disminuir el límite de
        detección aumenta el número de estrellas detectadas.
\end{itemize}

\end{document}